\cotefile{15}
\scanpage{170}

\uline{Foncteurs fibres motiviques, et \unsure{gpes} \unsure{fondamentaux} motiviques.}

À examiner :

0) VA \textcircled{à la Weil} \unsure{associées} aux motifs purs de poids impair \unsure{sous}
\unsure{Conditions} de Hodge--Riemann\ldots

\margin{OK *}
\textcircled{1} Extension des \unsure{groupes} \unsure{définis} par des \unsure{présentiques|présentations}
\uline{\unsure{Groupe} fondamental motivique.}

\textcircled{2} Relations avec le corps de classes

\margin{OK}
\textcircled{3} Cas du spectre d'un corps fini. Pb motivique
\unsure{points|partiel} à \unsure{val.} dans $\mathbb{R}$ ???

\margin{*}
4) Rôle des \unsure{gpes} fondamentaux motiviques pour la comparaison des réseaux canoniques
\diagram{\[
\left[
\begin{array}{ll}
\text{des Coh. de Hodge} &
\left|\begin{array}{l}
\text{2 réseaux \unsure{transcendts|transcendants} \unsure{pr.|pour} des places diff}\vphantom{\big|}\\
\text{réseau transcendant et réseau de Hodge}
\end{array}\right.\\[6pt]
\text{des Coh. $l$-adique} &\big|\ \text{2 réseaux transcendants}
\end{array}
\right]
\]}

5) Relations \inserted{de 4)} avec mes notes anciennes

6) Fonctions $L$ : équation fonctionnelle

\margin{OK}
7) Structures \unsure{positives}\ldots

\deleted{\unsure{7}}8) Topologie à mettre sur $\mathbb{Z}$ $-\, -\, -\, -$

9) Hyp. de Riemann ??

\margin{\unsure{Positivité} à la Hodge en se ramenant à des questions de \uline{diviseurs} par \unsure{éclatement|éclatements}, ou \unsure{vertu des|voie des} \unsure{pinceaux} ??}

\margin{169}

\scanpage{171}

SYSTÈMES de RACINES (20)\hfill Tribu n\textsuperscript{o} 59

de signe $+1$.

1) On pose $E=\unsure{\Pi|Q}/3P$. C'est un $\uuline{F}_3$-espace vectoriel de dimension 5. Montrer que le produit $(x,y)$ définit sur $E$ une forme bil. sym. non dégénérée. Montrer que deux éléments distincts de $R$ ont des images distinctes dans $E$.

2) On note $O(E)$ le groupe orthogonal de $E$, muni de la forme bil. précédente. On a $O(E)=\{\pm1\}\times SO(E)$. La norme spinorielle définit un homomorphisme surjectif $SO(E)\longrightarrow\{\pm1\}$ dont le noyau est noté $SO^+(E)$. Le groupe $SO^+(E)$ est simple, (d'ordre $25\ 920$) ; le quotient $O(E)/SO^+(E)$ est de type $(2,2)$. En déduire que $O(E)$ contient un sous-groupe $\Omega(E)$ d'indice 2 distinct de $SO(E)$ et ne contenant pas $-1$ ; montrer que ce sous-groupe est unique.

3) Tout élément de $\operatorname{Aut}(R)$ définit par passage au quotient un élément de $O(E)$. Montrer, (en comparant les ordres), que l'on obtient ainsi un isomorphisme de $\operatorname{Aut}(R)$ sur $O(E)$ ; l'image de $W(R)$ par cet isomorphisme est $\Omega(E)$, celle de $W^+(R)$ est $SO^+(E)$.

(Ainsi, $W(R)$ est extension d'un groupe à deux éléments par le groupe simple d'ordre $25\ 920$).

F) Groupe de Weyl $W$ de $E_7$.

Soient $Q$ le réseau radiciel, $P$ celui des poids. On a $2P\subset Q$. Posons $E=Q/2P$. La forme bilinéaire $(x\ y)$ définit, par passage au quotient, une forme bilinéaire sur le $F_2$-espace vectoriel $E$, (noter que $(Q\mid 2P)\subset 2Z$). Montrer que cette forme est non dégénérée et alternée, (noter que $\dim(E)=6$). En déduire un homomorphisme $W\longrightarrow\operatorname{Sp}(E)=\operatorname{Sp}(3,F_2)$.

Montrer que la suite
\diagram{\[
\text{(s)}\qquad
\begin{tikzcd}[column sep=large]
1 \arrow[r] & (1,-1) \arrow[r] & W \arrow[r] & \operatorname{Sp}(3,F_2) \arrow[r] & 1
\end{tikzcd}
\]}
est exacte, (on peut utiliser le fait que $\operatorname{Sp}(3,F_2)$ est d'ordre $2^9.3^4.5.7$ ; cf. Alg. IX).

Montrer que le sous-groupe $W^+$ formé des éléments de déterminant 1 de $W$ se projette isomorphiquement sur $\operatorname{Sp}(3,F_2)$, (noter que ce sous-groupe ne contient pas $-1$).

Donner une seconde démonstration de (s) en utilisant la structure de $W(E_8)$ et l'isomorphisme $\operatorname{Sp}(3,F_2)\longrightarrow\underline{O}(7,F_2)$.

\margin{170}

\scanpage{179}

\deleted{On peut supposer $X$ irréductible de pt générique $\eta$.
L'inclusion $\unsure{\eta-X|\eta\to X}$ définit une suite spectrale
(pour la cohomologie étale)
\[
H^*(\eta,G_m)\Leftarrow E_2^{pq}=H^p(X,R^qi_*(G_m))
\]
Or \deleted{\unsure{comme les}} comme les $H^q(\eta,G_m)$ et les fibres de $R^qi_*(G_m)$, qui s'explicitent par de la cohomologie galoisienne, sont des groupes de torsion pour $q>0$,
\unsure{je trouve que} \deleted{\unsure{il en va de même des} \unsure{gpes|groupes}}
\[
(*)\qquad H^p(X,R_X^*)\quad
\deleted{,\quad R_X^*=R\unsure{i_*}(G_m)}
\]
\inserted{est un groupe de torsion pour $p>0$,}
où $R_X^*=\deleted{\unsure{Ri_*}}i_*(G_m)$ est le faisceau (étale) des fonctions rationnelles \inserted{non nulles} sur $X$. \deleted{(\unsure{ceci étant}}}

D'autre part, on a \inserted{si $X$ est \inserted{noethérien} ``sans idéaux premiers immergés'',} la suite exacte
\diagram{\[
\text{(1.1.)}\qquad
\begin{tikzcd}[column sep=large]
0 \arrow[r,dash] & G_{mX} \arrow[r] & R_X^* \arrow[r] & \underline{\operatorname{Div}}_X \arrow[r] & 0
\end{tikzcd}
\]}
où $\underline{\operatorname{Div}}_X$ est le faisceau (\inserted{pour la top.} étale) des diviseurs (de Cartier) sur $X$, ce qui \deleted{\unsure{veut dire que}} \inserted{implique} compte tenu de ce que
\diagram{\[
\begin{tikzcd}[column sep=large]
H^{\unsure{i}}(X,\underline{\operatorname{Div}}_X) \arrow[r,"\partial"] & H^{i+1}(X,G_{mX})
\end{tikzcd}
\]}
\margin{\uline{\unsure{donne le|donner le}} \uline{corollaire} \unsure{3.4.}}
est bijectif modulo torsion \deleted{pour} si $i\geq 1$. \uline{Ainsi}
\deleted{\unsure{pour en montrer l'utilité}} L'hypothèse de factorialité d'autre part implique \deleted{l'isomorphisme} qu'on a un isomorphisme de faisceaux étales :

\margin{178}

\scanpage{180}

\unsure{Calcul|Calculs} de \unsure{cohomol.} absolue

du motif de Tate
\[
\mathbb{Z}(\unsure{n|u}).
\]

\margin{179}

\scanpage{181}
[Redacted: a private administrative document on the back of this sheet, not mathematics.]

\scanpage{182}

Généralités \unsure{fonctoriello-catégoriques}
sur la distribution des
frobenius dans les \unsure{gpes}
de Galois motiviques.

\margin{181}

\scanpage{183}

Soit $\mathcal M$ une catégorie de motifs \unsure{géomologiques|géométriques}, \inserted{de type fini,} sur le préschéma $S$ de type fini sur $\mathbb Z$. Pour simplifier, supposons que les degrés de $\mathcal M$ soient pairs (sinon, il sera possible, pour \deleted{\unsure{la plupart des}} \inserted{certaines} questions de distribution de Frobenius, de remplacer $\mathcal M$ par $\mathcal M^{\unsure{\mathrm{pairs}|\mathrm{pair}}}$). Tout choix d'un pt fermé $s\in S$ permet \unsure{donc}, de définir des (\unsure{à iso bien déf.}) ``$\mathbb R$-points motiviques'' \deleted{\unsure{de $\mathcal M$}''}\inserted{de $\mathcal M$}.

Pour fixer les idées, choisissons un pt $\mathbb R$-motivique de $\mathcal M$, d'où un gpe algébrique $G_{\mathbb R}$ sur $\mathbb R$.

Il y a lieu de \deleted{\uline{\unsure{dédoubler}}} \inserted{démultiplier} les points \inserted{fermés $s$} de $S$, en prenant les couples $s':(s,i)$, où $s\in S^0$, et
\margin{\unsure{l'exposant} $^0$ désigne \unsure{les} pts \uline{fermés}}
$i:\unsure{G_{m\mathbb C}-G^{(s)}_{\mathbb C}|G_{m\mathbb C}\to G^{(s)}_{\mathbb C}}$ est de type \deleted{\unsure{donné}} préférentiel -- $G(s)$ désigne le gpe de Galois motivique associé à $\mathcal M_s$.

\deleted{Alors} Soit $S'$ l'ens de ces couples. À tout $s'\in S'$ est associé un $\mathbb R$-point motivique de $\mathcal M$, \deleted{d'où un} déterminé à isom près, d'où un
\[
\unsure{\xi|\zeta}(s')\in H^1(\mathbb R,\unsure{G^{00}_{\mathbb R}}),
\]
où $G^{\unsure{00}}$ est le noyau de
\diagram{\[
\begin{tikzcd}
G^0_{\mathbb R} \arrow[r,"\varepsilon"] & G_m .
\end{tikzcd}
\]}
De plus, si $G^{\unsure{\xi}(s')}_{\mathbb R}$ est ``le'' gpe déduit de $G_{\mathbb R}$ par torsion par $\unsure{\xi}(s')$, on obtient un hom
\diagram{\[
\begin{tikzcd}[column sep=large]
u(s'):\ G(s')_{\deleted{\unsure{\mathbb R}}\mathbb R}=G(s)_{\mathbb R} \arrow[r] & G^{\unsure{\xi}(s')}_{\mathbb R}
\end{tikzcd}
\]}
compatible aux $\varepsilon_{\mathbb R},\unsure{J_{\mathbb R}|j_{\mathbb R}}$, et déterminé en fait mod \deleted{automorphisme intérieur} \inserted{conjugaison} par un élément de $G^{\unsure{00}}_{\mathbb R}(\mathbb R)$.

\margin{182}

\scanpage{184}
\mspage{73}

a homomorphism of restriction of a function on $\mathcal E_{\varphi}$. We obtain from (16):
\diagram{\[
\begin{tikzcd}[column sep=large]
0 \arrow[r] & O_V(-\mathcal E_{\varphi})[nE] \arrow[r] & O_V[nE] \arrow[r,"r'"] & O_{\mathcal E_{\varphi}}[nE]_{\mathcal E_{\varphi}} \arrow[r] & 0
\end{tikzcd}
\qquad(17)
\]}

here $r'$ is such: if $P\notin\mathcal E_{\varphi}$, then it is a zero homomorphism; if $P\in\mathcal E_{\varphi}$ and $\{\alpha^{(i)}\}\in O_V[nE](P)$, then $r'\{\alpha^{(i)}\}=\{r\alpha^{(i)}\}$. To the sequence (17) corresponds such exact sequence:
\diagram{\begin{gather*}
\begin{tikzcd}[column sep=large]
H^0(V,O_V[nE]) \arrow[r,"r^*"] & H^0(\mathcal E_{\varphi},O_{\mathcal E_{\varphi}}[nE]_{\mathcal E_{\varphi}}) \arrow[r] & H^1(V,O_V(-\mathcal E_{\varphi})[nE]) \arrow[r] & {}
\end{tikzcd}\\
\begin{tikzcd}[column sep=large]
{} \arrow[r] & H^1(V,O_V[nE]) \arrow[r] & H^1(\mathcal E_{\varphi},O_{\mathcal E_{\varphi}}[nE]_{\mathcal E_{\varphi}})
\end{tikzcd}\qquad(18)
\end{gather*}}

Let $n(\mathcal E_{\varphi})$ be such that, at $n\geq n(\mathcal E_{\varphi})$ for $\mathcal E_{\varphi}$ and $[nE]_{\mathcal E_{\varphi}}$ satisfy the requirements 1') and 2') (See p. 202). The existence of such $n(\mathcal E_{\varphi})$ follows from the condition of non-degeneracy of $[E]_{\mathcal E_{\varphi}}$ over $\mathcal E_{\varphi}$. Put
\[
n_{\mathcal E}=\max(n'',n(\mathcal E_{\varphi})).
\]

Then, at $n\geq n_{\mathcal E}$, from (15), (18) and from the fact that, at $n\geq n(\mathcal E_{\varphi})$

\margin{183}

\scanpage{185}

De plus, on a une structure \unsure{polarisée} \inserted{$\pi$} sur $G^0_{\mathbb R}$, qui est \deleted{\unsure{déf}} supposée définie par un hom
$\deleted{\unsure{\mu_{2\mathbb C}}}\quad i:\mu_{2\mathbb C}\unsure{-|\to}G^0_{\mathbb C}$
(comme on voit en envoyant un \unsure{tore} \inserted{point d'ordre 2} dedans). D'autre part, $G(s')_{\mathbb C}$ est muni de $i(s'):G_{m\mathbb C}\to G^{(s')}_{\mathbb C}$, et on obtient par conséquent
\diagram{\[
\begin{tikzcd}[column sep=large]
\unsure{v|v_q}(s')=u(s')_{\mathbb C}\circ(i(s')\mid\mu_2):\ \mu_{2\mathbb C} \arrow[r] & G^{\unsure{\xi}(s')}_{\mathbb C}
\end{tikzcd}
\]}

\deleted{D'autre part} Ceci posé, il faut que $\pi^{\unsure{\xi}(s')}$ (polarisation de $G^{\unsure{\xi}(s')}_{\mathbb R}$ déduite de $\pi$ sur $\mathbb C$ par transport) soit celle définie par les $\mu_{2\mathbb C}$.

Par conséquent, pour $s'$ \unsure{variable|variant} de façon que $\unsure{\xi}(s')$ soit constant, les polarisations \inserted{de $G^{\unsure{\xi}(s')}_{\mathbb R}$} associées aux différents $v(s')$ doivent être identiques -- ce qui donne une condition non tautologique sur la distribution des \deleted{\unsure{ces}} $u(s')$. [Noter cependant que l'ens des \deleted{classes de} polarisations d'un gpe réel \inserted{$H$}, définies en termes d'un hom $\mu_{\mathbb C}\to H_{\mathbb C}$, est \uline{fini}, [si $G$ est commutatif en tous cas !], donc la restriction n'est pas tellement draconienne.
\margin{? à vérifier}

Enfin, conformément aux conjectures de Tate \unsure{magnifiées}, La connaissance des $u(s')$ permet de retrouver la structure d'effectivité dans la catégorie

\margin{184}

\scanpage{186}
\mspage{72}

Since
\[
O_V(-\unsure{E_{\varphi}|\mathcal E_{\varphi}})[nE]\cong O_V[(n-1)E]
\]
then, having taken $n''\geq n'$, we obtain from (14) that, at $n\geq n''$
\[
\begin{gathered}
\dim H^1(V,O_V[nE])=\\[6pt]
=\dim H^1(V,O_V(-\mathcal E_{\varphi})[nE])=\dim H^1(V,O_V[(n-1)E]).
\end{gathered}\qquad(15)
\]

\margin{p. 208}
It is evident that $\varphi^n=\{\varphi_i^n\}$ is a cross-section of the group $H^0(V,O_V[nE])$. Therefore, the common zeros of the elements of the group $H^0(V,O_V[nE])$ are contained among the zeros (of) $\varphi^n$, that is, all of them lie on the subvariety $\mathcal E_{\varphi}$.

Let us show that, at $n\geq n_{\mathcal E}$ there is no common zeros of the cross-sections of $H^0(V,O_V[nE])$ on $\mathcal E_{\varphi}$.
Consider an exact sequence
\diagram{\[
\begin{tikzcd}[column sep=large]
0 \arrow[r] & O_V(-\mathcal E_{\varphi}) \arrow[r] & O_V \arrow[r,"r"] & O_{\mathcal E_{\varphi}} \arrow[r] & 0
\end{tikzcd}\qquad(16)
\]}

The homomorphism $r$ is described in the following way: if $P\notin\mathcal E_{\varphi}$, then $r$ in the stalk over this point is a zero homomorphism; if $P\in\mathcal E_{\varphi}$, then $r$ over $P$ is

\margin{185}

\scanpage{187}
des $G_{\mathbb R}$-modules de dim finie sur $\mathbb R$. Le cône
qui exprime cette effectivité est \uline{engendré}
par les \deleted{$v(s')$} rayons définis par les $v(s')$.
Cela ne fait qu'exprimer \inserted{sous forme ``analytique''} le fait que les
frobenius $f(s')\in G_{\mathbb R}^{\unsure{\xi}(s')}(\mathbb R)$ sont tels
que leurs valeurs propres, dans tte représentation
de $G_{\mathbb R}\deleted{(\unsure{\mathbb R})}^{\unsure{\xi}(s')}$ --- et en fait $\unsure{\overline{\mathrm m}|\simeq}$ de
$G_{\deleted{\mathbb R}\inserted{\mathbb C}}^{\unsure{\xi}(s')}$ --- sont des \uline{entiers} algébriques, ce qui est une
condition arithmétique assez draconienne !
\margin{186}

\scanpage{188}
\mspage{70}

that there exists $n_0(K)$ such that at $n\unsure{\mathrm m|\geq}n_0(K)$
\[
H^i(V,K[nE])=0\qquad i=1,2,\ldots
\]

From the exact sequence
\diagram{\[
\begin{tikzcd}[column sep=large]
H^i(V,O_V(-E_{\varphi})[nE]) \arrow[r] & H^i(V,O_V[nE]) \arrow[r] & H^i(V,K[nE])
\end{tikzcd}
\]}
at $n\geq n_0(K)$, $i=1,2,\ldots$ we obtain:
\[
\dim H^i(V,O_V[nE])\leq\dim H^i(V,O_V(-E_{\varphi})[nE]).
\]

Since
\[
O_V(-E_{\varphi})\cong O_V[-E],
\]
then
\[
O_V(-E_{\varphi})[nE]\cong O_V[-E][nE]\cong O_V[(n-1)E].
\]

Therefore, at $n\geq n_0(K)$, $i=1,2,\ldots$
\[
\dim H^i(V,O_V[nE])\leq\dim H^i(V,O_V[[n-1]E]).
\]

Evidently, there exists such $n'$ that, at $n\geq n'$

\margin{187}

\scanpage{189}
\uline{Théorie de Galois des Motifs}

\uline{Généralités}

\margin{188}

\scanpage{190}
Foncteurs fibres motiviques
et \unsure{gpe|gpes} \unsure{fondamental|fondamentaux}
\unsure{motivique|motiviques}

\margin{189}

\scanpage{191}
$X$ préschém. connexe \deleted{\unsure{finitude} \unsure{a priori} \unsure{quelconque}}

$1^\circ)$ $X_{\mathrm{réd}}$ de car $p>0$. Alors dans le pro-torseur fondamental
$(T_\alpha)$, il n'y a pas en général de sections pour les $T_\alpha$.

\margin{$\unsure{\ast}$}
On dispose uniquement, \deleted{de façon canonique}

\margin{Coh. $l$-adique}
(i) des $\unsure{\xi}_l^{(\alpha)}\in T_\alpha(\unsure{\mathbb Z|\mathbb Q}_l)$,
$l\neq p$, provenant sans doute
de $\unsure{\xi}_l\in T_\alpha(\widetilde{\mathbb Z}_{l\mathbb Z})$.
\margin{hensélisé}

\margin{\fbox{\begin{minipage}{.29\linewidth}
Coh. $p$-adique\\
Structure supplémentaire :\\
opérations \fbox{$\uuline{F}\text{ et }\uuline{V}$}\\
$1/2$ \uuline{linéaires} \uuline{\unsure{relatifs}}\\
\unsure{Frobenius sur $k$}
\end{minipage}}}

\diagram{\[
\begin{tikzcd}[column sep=small,row sep=large]
\text{(ii) d'un }\unsure{\xi}_p^{(\alpha)}\in T_\alpha( &
\mathbb W_p(k) \arrow[dr,"\text{``DeRham-Hodge''}"'] &
\begin{gathered}
)\quad\text{si }\deleted{\text{$X$ déf. sur}}\inserted{\text{$k$ est un corps parfait sur $X$}}\\
\deleted{\text{le corps parfait $k$}}\\
[\ \unsure{\xi}_p\ \text{provenant sans doute}\\
\text{d'un }\unsure{\xi}_p\in T_{\unsure{p|\alpha}}(\widetilde{\mathbb Z}_{p\mathbb Z})\ ]
\end{gathered}
\\
\text{(iii)}[\ ?\quad\text{\uline{réduit}}\ \unsure{\xi}_p^{(\alpha)\unsure{0|\circ}}\in T_\alpha( &
\deleted{\unsure{\mathbb W}} & k)\ ]
\end{tikzcd}
\]}
\margin{hensélisé strict.}

\margin{Coh. Hodge-deRham. Mais ici, il faut \deleted{sans doute}
travailler dans la catégorie dérivée. Tenir compte aussi de la
structure \uline{filtrée} de Hodge-DeRham}

$\overline{\mathrm M}$ si $X=\operatorname{Spec}\mathbb F_q$, il y a en général
des \deleted{$\unsure{\xi}$} \unsure{obstructions} \unsure{non triviales} dans
des $H^{\unsure{2|1}}(\mathbb Q,G)$ $-\, -\, -$

\margin{$\unsure{\ast}$}
On utilise des ``foncteurs fibres'' qui induisent sur
les motifs effectifs de poids $0$ le foncteur
canonique. Alors les objets de $T_\alpha(\mathbb Z_l)$ sont
sans doute uniques\inserted{à isom près}, (th. de Lang), de $\unsure{\overline{\mathrm m}|\simeq}$
ceux de $T_\alpha(\unsure{\mathbb W}_p(k))$ si $k$ fini (ib), \inserted{plus gén.} pour un
$k$ parfait de dim \unsure{ch} $1$ (?). En général, il
faut s'attendre que $T_\alpha(k)$ \deleted{soit}\inserted{ait} plus d'un
\unsure{élém.} à isom près, mais il est possible
que $T_{\unsure{\infty}}(k)=\varprojlim T_\alpha(k)$ n'a qu'un \unsure{élém.} à
isom près, celui donné par la \unsure{théorie|théor.} de DeRham
Hodge.

\margin{\unsure{relation avec|relation au} \unsure{le problème} \unsure{galoisien} ??
Il doit être \uline{\unsure{épuisant|équivalent}}, cf \uline{Manin}}

\margin{190}

\scanpage{192}
ALGÈBRES de LIE SEMI-SIMPLES (9)\hfill Tribu n$^\circ$ 59

Pas de noms propres à la p. 55 ; Chevalley et Harish on parlera longuement
de dans la note historique. Dans la prop. 7 appeler les gens donnés $X'$ et $Y'$.
Au lieu du ``tout de suite'', (l.2, p. 56), dire pourquoi, (ça se propage à partir
des racines simples).

Utiliser le lemme sur les applications $R\longrightarrow$ groupe au milieu de la p. 56.
Ensuite poser $X_\alpha=u_\alpha^{-1}X'_\alpha$, (alors $t_\alpha=1$). Pas de nom à la p. 57.

\uline{\unsure{Exer6|Exerb}} - Dans l'exer. 3 donner le système des $h$, $x$, $y$ de Weyl, qui fournit des
formules sans dénominateurs. Donner la démonstration Van der Waerden, (cf.
Jacobson, bas p. 123 et p. 124).

\S4 - \uuline{Théorèmes d'existence}, (nouveau titre).

Dans ce \S\ on fait deux constructions distinctes, qui amènent presque aux
mêmes résultats. On en laissera tomber une, d'où disparition des th. 3 et 4.
Ceci amène au plan suivant :

1) Construction de $\underline a=\underline x+\underline h+\underline y$ ;

2) Définition des gros modules $E(\lambda)$ ; liberté de $\underline x$ et $\underline y$ ;

3) Construction de $\underline g$ comme quotient de $\underline a$ par les $\unsure{\Theta|\theta}_{ij}$ ; alors
$\underline g=\underline h+\sum_\lambda\underline g^\lambda$,
est de dim. finie, et est semi-simple ;

4) Définition des $W(\lambda)$ ; sont des $g$-modules ;

5) $W(\lambda)$ est de dim. finie si $\lambda$ est un pouah dominant ; conséquences pour $\underline g$.

Dans 3) la semi-simplicité de $\underline g$ se démontre ainsi : on a
$[\underline g^\alpha,\underline g^{-\alpha}]=\underline h^\alpha$,
(vrai pour $\alpha$ simple, et $W$ opère) ; alors tout idéal commutatif $\underline b$ est nul, (car
$\underline b\cap\underline g_\alpha=0$ pour toute sous-algèbre $\underline g_\alpha$ de dim. 3). A partir de là tous les cor.
concernant $\underline g$ se démontrent facilement.

\margin{191}

\scanpage{193}
Les $\unsure{\xi}_l$ de (i) dépendent du choix d'un pt géom de $X$,
mais à isom non unique près (dépendant du
\uuline{$\pi_1(X)$}) \inserted{[$\pi_1(X)$ opère sur $\unsure{\xi}_l$]}
\deleted{sont} \unsure{indépendants} de ce choix.

Les $\unsure{\xi}_p^\alpha$ dépendent du choix d'un \deleted{corps} pt\inserted{de $X$}
à valeurs dans un corps parfait sans plus,
si on veut les $k(x)^{\mathrm{pf}}$ (clôture parfaite de $k(x)$),
et pour $k$ donné, il ne \uline{semble pas} qu'il y ait
\deleted{ait} des isom entre les $\unsure{\xi}_{p,a}^{(\alpha)}$ correspondants
aux divers pts à valeurs dans $k$ (c'est la
``partie continue'' du $\unsure{H^1|H^2}$ préc. qui v. faire des
ennuis), s'ils sont localisés en le $\overline{\mathrm m}$ $x\in X$, \deleted{il}\inserted{et}
\deleted{sont} si $k$ est alg. clos, ils sont conjugués par
autom. de $k$\inserted{(étendus à $\mathbb W_p(k)$)}, les groupes alg. associés sont
donc conjugués par autom. de $k$ $-\, -\, -\, -$ il est $\overline{\mathrm m}$
gentil qu'ils \unsure{veulent|veuillent} être l'un \unsure{tordu} de
l'autre $-\, -\, -\, -$ Le gpe $\pi_1(X)$ n'opère pas sur $\unsure{\xi}_p^\alpha$
\figure{\begin{tikzpicture}[baseline=-.4ex,scale=.22]\draw (0,0)--(1,2);\draw (1,2) circle (.18);\end{tikzpicture}}, par
contre le groupe des autom de $k$ sur $k(x)$
\deleted{[\unsure{au $\overline{\mathrm m}$ de $k$} \unsure{$H$ précédent}]}
opère de façon $1/2$ \uline{linéaire} sur $\unsure{\xi}_p^{(\alpha)}$ $-\, -\, -\, -$
\margin{[\unsure{En une descente}
\figure{\begin{tikzpicture}[baseline=-.4ex,scale=.22]\draw (0,0)--(1,2);\draw (0,0) circle (.18);\end{tikzpicture}}}

\margin{\figure{\begin{tikzpicture}[baseline=-.4ex,scale=.22]\draw (0,0)--(1,2);\draw (1,2) circle (.18);\end{tikzpicture}}
En revanche, on a comme structure supplémentaire
les opérations $\unsure{F\,V|F,V}$ avec $FV=p$}

\margin{\figure{\begin{tikzpicture}[baseline=-.4ex,scale=.22]\draw (0,0)--(1,2);\draw (0,0) circle (.18);\end{tikzpicture}}
\unsure{élégante.} D'autre part, un autom $\sigma$
de $k$ qui n'est pas $k(x)$-linéaire
transforme en général $\unsure{\xi}_{xp}^{\unsure{\alpha}}$ en un
$\unsure{\xi}_{xp}^{\sigma}$ qui n'est plus isom
à $\unsure{\xi}_{x,p}$ $-\, -\, -\, -$}

Notations

\diagram{\[
\text{(i)}\left\{\begin{array}{ll}
\unsure{\xi}_{\bar x,l} & \text{foncteurs fibres $l$-adiques en}\\
& \leftarrow\ \text{pt géom }\bar x\\[6pt]
\unsure{\xi}_{x,l} & \text{par desc. de la \unsure{$2$-gpe|$l$-gpe} pour $x\in X$}\\
& \text{ou un pt de $X$ à val.}\\
& \text{dans un corps quelc.}\\
& \text{($\unsure{\xi}_{x,l}$ déf. à isom \uline{non uniq.} près)}
\end{array}\right.
\]}

\margin{\uline{On peut mettre}\\
\uline{des \unsure{exposants} $\alpha$}\\
\uline{à \unsure{volonté}}}

\diagram{\[
\text{(ii)}\left\{\begin{array}{ll}
\unsure{\xi}_{x,p} & \text{$x$ \unsure{est pt} de $X$ à val. dans un corps}\\[6pt]
\unsure{\xi}_{x,p}^{0} &
\end{array}\right.
\]}

\margin{192}

\scanpage{194}
ALGÈBRES de LIE SEMI-SIMPLES (8)\hfill Tribu n$^\circ$ 59

Dire ``ces éléments $N_{\alpha\beta}$'' à la p. 46, l.6 bas.

\uline{n$^\circ$ 6} - Enoncer ainsi la prop. 2 : a) tout sous-espace $V$ stable par $\underline h$ est de la
forme, (sous-espace $\underline h'$ de $\underline h$) $\oplus\displaystyle\sum_{\alpha\in\Sigma}\underline g^\alpha$ ; b) pour que $V$ soit une sous-algèbre,
il faut et il suffit que $\Sigma$ soit clos et que $\alpha\in\Sigma\cap(-\Sigma)$ implique $H_\alpha\in\underline h'$.
Exhiber alors le radical de $V$, (une partie vient de $\underline h$, une autre des $\alpha\in\Sigma$
tels que $-\alpha\notin\Sigma$).

Enoncé de la prop. 3 : $V$ réductive $\Longleftrightarrow$ $\Sigma$ clos et symétrique ; $V$ semi-simple $\Longleftrightarrow$
$\Sigma$ clos et symétrique et $\underline h'$ engendré par les $H_\alpha$ pour $\alpha\in\Sigma$ ; son
système de racines est alors celui que dit la prop. 3.

Le cor. 2, p. 47, ira en \unsure{exer6|exerb}. La dém. de la prop. 4 se simplifie, maintenant
qu'on connaît le radical. Dans le cor., p. 51, dire ``racines positives par
rapport à une base de $R$''.

A propos des Borels, (p. 51 bas), dire que la forme quadratique est nulle
sur $\underline n^+$ et $\underline n^-$, et que leurs orthogonaux sont $\underline b^-$ et $\underline b^+$.
On parlera des paraboliques, (ou paraboréliques).

\uline{n$^\circ$ 7} - Dans la prop. 7, indexer $S$, (qui s'appelle maintenant ``base de $R$'').
A la l.8 bas, (p. 52), dire que les autres formules ``ont déjà été démontrées'',
et donner les références.

Donner ici les relations $(\operatorname{ad}X_i)^{-a_{ij}+1}X_j=(\operatorname{ad}Y_i)^{-a_{ij}+1}Y_j=0$, (cf. \S\ 4).

Pas de nom pour $(H_i,\,X_i,\,Y_i)$ en haut de la p. 53, ça conflicte ; on dira
que ce système est associé à l'épinglage.

\uline{n$^\circ$ 8} - Au début préciser que $X_\alpha\in\underline g^\alpha$. Le lemme 4 remonte au \S\ 2, n$^\circ$ 2. Introduire
$p$ aussi dans l'énoncé de la prop. 6. La formule du bas de la p. 54 tient
à ce que $\alpha+\beta$ est une racine ; le dire.

\margin{193}
